% TOP_PROBLEM: {"id": 3270, "title": "Hoàng-Reed Conjecture", "category_id": 3, "category": {"id": 3, "name": "graph_theory", "display_name": "Graph Theory", "description": "Problems involving graphs, networks, and their properties.", "slug": "graph-theory", "order_index": 3, "created_at": "2026-07-31T15:26:25.671Z"}, "status": "open", "problem_number": "OPG-47282", "set_id": 12}
% FIRST_POSTED: 2026-10-01
% ENTRY_KIND: statement_only
% UnsolvedMath ID 3270; source catalogue code OPG-47282
% !TeX program = lualatex
% Definition and sources only; this file is not an LLM solution attempt.
\documentclass[11pt,a4paper]{article}
\usepackage[margin=25mm]{geometry}
\usepackage{fontspec}
\setmainfont{lmroman10-regular.otf}[BoldFont=lmroman10-bold.otf,
 ItalicFont=lmroman10-italic.otf,BoldItalicFont=lmroman10-bolditalic.otf]
\usepackage{amsmath,amssymb,mathtools}
\usepackage{unicode-math}
\setmathfont{latinmodern-math.otf}
\AtBeginDocument{\let\setminus\smallsetminus}
\usepackage{xurl,hyperref}
\hypersetup{colorlinks=true,urlcolor=blue}
\setcounter{secnumdepth}{0}
\setlength{\parindent}{0pt}
\setlength{\parskip}{5pt}
\setlength{\emergencystretch}{3em}
\begin{document}
\section*{Hoàng-Reed Conjecture}\label{3270}
{\small UnsolvedMath ID 3270 · OPG-47282 · Graph Theory · OpenGarden}
\subsection{Definitions and mathematical statement}
Let $D=(V,A)$ be a nonempty finite loopless simple digraph, allowing opposite
arcs. Write $d^+(v)$ for the number of distinct vertices $w$ with
$vw\in A$. A directed cycle is a subdigraph following a cyclic sequence
of distinct vertices; cycles of length two are allowed when both
opposite arcs occur.

For each integer $k\ge1$, the Hoàng--Reed conjecture asks whether the
condition $d^+(v)\ge k$ for every vertex guarantees directed cycles
$C_1,\ldots,C_k$ which can be ordered so that
\[
 \left|V(C_j)\cap\bigcup_{i=1}^{j-1}V(C_i)\right|\le1
                         \qquad(2\le j\le k).
\]
Thus each newly chosen cycle is either disjoint from all earlier cycles
or meets their entire union at just one vertex. In particular the
cycles are arc-disjoint. They are not required to be vertex-disjoint,
and an arbitrarily large number of them may use one common vertex.

The condition is stronger than saying that every pair of cycles shares
at most one vertex: a new cycle cannot meet two earlier cycles at two
different vertices. The existence of a suitable ordering is part of
the conclusion. Extra arcs of $D$ between vertices of the chosen cycles
are ignored, since the union is not required to be induced.

For $k=1$ the assertion reduces to finding a directed cycle in a finite
digraph with no sink. The general question seeks a structured collection
of $k$ cycles from a local outdegree assumption, without any minimum
indegree or connectivity requirement.
\subsection{Short English statement}
Does minimum outdegree k force k directed cycles that can be added one by one, each meeting all previous cycles in at most one vertex?
\subsection{Sources}
\begin{itemize}
\item[S1] ULAM.AI and the UnsolvedMath Contributors, supplied problems.json, record 3270 (OPG-47282), OpenGarden collection. Catalogue identity and source transcription; CC BY 4.0.
\url{https://huggingface.co/datasets/ulamai/UnsolvedMath}.
\item[S2] Open Problem Garden, Hoàng-Reed Conjecture. Original problem and discussion; the mathematical formulation below is expanded from the supplied source record.
\url{http://www.openproblemgarden.org/op/hoand_reed_conjecture}.
\end{itemize}
\end{document}
